\chapter{C++}\label{ch:iv:cpp}

``What does this print?'' Eighteen lines: a base class whose constructor calls a virtual function, a derived class that
overrides it, one object. The candidate answers ``Derived'', explains dynamic dispatch well, and is wrong: while the base's
constructor runs, the object is still a base, and the base's version is called. C++ interviews at trading firms test the
language as the machine sees it: when objects begin and end, what a move really moves, what the compiler may assume, and
what the standard library does inside. The performance side is One Quant Book~13's (chapters 6 to~8); this chapter is the
interview side. Every snippet printed here is compiled and run by the chapter's tests, and every claim of undefined behaviour
is diagnosed by a sanitizer, never ``tested'' for an output.

\section{Object lifetime}

\begin{definition}[Output-prediction question]\label{def:iv:cpp:output}
An \emph{output-prediction question}\index{output-prediction question} shows a short program and asks what it prints. It
tests the language's rules (order of construction and destruction, overload resolution, value categories, conversions)
and the candidate's ability to say when the answer is guaranteed, implementation-defined, unspecified, or not defined at all.
\end{definition}

The rules of lifetime that such questions probe are few. Members are constructed in the order they are declared, not the
order of the initialiser list, and destroyed in reverse. A temporary lives until the end of the full expression that created
it, unless it is bound directly to a local reference, which extends its life. During a constructor or destructor the dynamic
type is the class being constructed or destroyed, so virtual calls do not reach derived overrides. Objects with static storage
in one translation unit are initialised in the order of their definitions; across translation units the order is unspecified.

\begin{method}[Answering an output-prediction question]\label{met:iv:cpp:output}
\begin{enumerate}
\item Walk the program statement by statement, tracking each object's birth and death.
\item At each call, apply overload resolution with the argument's value category, not its declared type.
\item Classify every doubtful line: guaranteed by the standard, implementation-defined (sizes, representations), unspecified
(the state of a moved-from standard container), or undefined (then say so and stop predicting).
\item Say what a compiler with warnings on would flag.
\end{enumerate}
\end{method}

\section{Moves, copies and value categories}

A class that owns a resource either owns it through members that manage themselves (a \texttt{std::vector}, a
\texttt{std::unique\_ptr}) and declares no special members at all, the rule of zero, or manages it directly and must define
all five: destructor, copy constructor, copy assignment, move constructor, move assignment. Moves should be \texttt{noexcept},
or standard containers will copy instead of moving when they grow. A moved-from standard object is valid but in an unspecified
state: it may be destroyed or assigned, and nothing else may be assumed. A named rvalue reference is itself an lvalue: inside a
function taking \texttt{std::string\&\& s}, passing \texttt{s} on selects the lvalue overload unless \texttt{std::move(s)} is
written.

\begin{example}[Which overload?]\label{ex:iv:cpp:overload}
\Cref{lst:iv:cpp:overload} calls \texttt{f} with a named string, a moved string, a temporary, and a named rvalue reference
inside \texttt{g}. The last one is the question's point: the parameter has a name, so it is an lvalue and the
\texttt{const std::string\&} overload is chosen.
\end{example}

\omcode{code/interviews/22-cpp/cpp/snippets/overload_rvalue.cpp}{1}{16}{Overload resolution by value
category.}\label{lst:iv:cpp:overload}

\section{Templates and compile-time code}

Templates are resolved at compile time and cost nothing at run time when used for static polymorphism (the curiously
recurring template pattern of One Quant Book~13, chapter~7). \texttt{constexpr} functions can run at compile time, and a
constant expression may not contain undefined behaviour, so the compiler rejects it: a table of powers of ten that multiplies
once too often fails to compile rather than overflowing silently (\cref{iq:iv:cpp:13}). Concepts state requirements on
template arguments in the signature, so that a wrong type fails at the call with a readable message.

\section{Undefined behaviour and the standard library's insides}

Undefined behaviour (One Quant Book~13, chapter~8) is not ``it crashes'': the compiler may assume it never happens and
optimise accordingly, so a program with it can do anything, including what the author expected until the next compiler
upgrade. The commonest cases in interviews: signed integer overflow, dangling references and views, iterators or references
invalidated by a container's reallocation, reading uninitialised values, and data races. The practical defence is the
sanitizers: AddressSanitizer finds use after free and out-of-bounds access, UndefinedBehaviorSanitizer finds overflow and
bad shifts, ThreadSanitizer finds races.

Three library facts recur. \texttt{std::vector} grows geometrically, so \texttt{push\_back} is amortised constant time and
reallocates occasionally, invalidating every pointer, reference and iterator into it; \texttt{reserve} avoids it when the size is
known. \texttt{std::unordered\_map} is a table of buckets with linked nodes, one allocation per element, which is why a
flat open-addressing table is faster on a hot path. \texttt{std::function} erases the type of any callable and stores it in a
small internal buffer or on the heap when it does not fit, which costs an allocation and an indirect call.

\section{Question bank}

\omcode{code/interviews/22-cpp/cpp/snippets/virtual_in_ctor.cpp}{1}{18}{A virtual call during
construction.}\label{lst:iv:cpp:virtual}

\begin{interviewq}[$\star$ \iqroles{developer} \iqfirm{market maker}]\label{iq:iv:cpp:1}
What does \cref{lst:iv:cpp:virtual} print, and why?
\end{interviewq}

\omcode{code/interviews/22-cpp/cpp/snippets/lifetime_order.cpp}{1}{21}{Members and a temporary.}\label{lst:iv:cpp:lifetime}

\begin{interviewq}[$\star$ \iqroles{developer} \iqfirm{proprietary firm}]\label{iq:iv:cpp:2}
What does \cref{lst:iv:cpp:lifetime} print, line by line?
\end{interviewq}

\begin{interviewq}[$\star$ \iqroles{developer} \iqfirm{any}]\label{iq:iv:cpp:3}
A class owns a raw array allocated with \texttt{new[]}. Which special member functions must it define, and how could you avoid
defining any?
\end{interviewq}

\begin{interviewq}[$\star$ \iqroles{developer} \iqfirm{proprietary firm}]\label{iq:iv:cpp:4}
What is RAII? Give two examples from a trading system.
\end{interviewq}

\omcode{code/interviews/22-cpp/cpp/snippets/smart_sizes.cpp}{1}{8}{Sizes of three library types.}\label{lst:iv:cpp:sizes}

\begin{interviewq}[$\star$ \iqroles{developer} \iqfirm{market maker}]\label{iq:iv:cpp:5}
Compare \texttt{std::unique\_ptr} and \texttt{std::shared\_ptr}: what does each cost, in memory and time? What does
\cref{lst:iv:cpp:sizes} print on x86-64 Linux with libstdc++, and which parts of that answer are guaranteed?
\end{interviewq}

\begin{interviewq}[$\star\star$ \iqroles{developer} \iqfirm{proprietary firm}]\label{iq:iv:cpp:6}
What does \cref{lst:iv:cpp:overload} print? Change one line so that \texttt{g} forwards its argument as an rvalue.
\end{interviewq}

\omcode{code/interviews/22-cpp/cpp/snippets/ub_dangling_view.cpp}{1}{10}{A view of a
symbol.}\label{lst:iv:cpp:view}

\begin{interviewq}[$\star\star$ \iqroles{developer} \iqfirm{market maker}]\label{iq:iv:cpp:7}
What is wrong with \cref{lst:iv:cpp:view}? How would you find the bug if you did not see it?
\end{interviewq}

\omcode{code/interviews/22-cpp/cpp/snippets/ub_signed_overflow.cpp}{1}{8}{Doubling a
quantity.}\label{lst:iv:cpp:overflow}

\begin{interviewq}[$\star\star$ \iqroles{developer} \iqfirm{proprietary firm}]\label{iq:iv:cpp:8}
What does \cref{lst:iv:cpp:overflow} print? What may an optimising compiler assume about a loop whose bound is \texttt{i + 1 >
i}?
\end{interviewq}

\omcode{code/interviews/22-cpp/cpp/snippets/ub_invalidated.cpp}{1}{10}{Keeping a reference to the first
fill.}\label{lst:iv:cpp:invalidated}

\begin{interviewq}[$\star\star$ \iqroles{developer} \iqfirm{market maker}]\label{iq:iv:cpp:9}
Find the bug in \cref{lst:iv:cpp:invalidated} and give two ways to fix it.
\end{interviewq}

\omcode{code/interviews/22-cpp/cpp/snippets/sign_compare.cpp}{1}{9}{A limit check.}\label{lst:iv:cpp:sign}

\begin{interviewq}[$\star\star$ \iqroles{developer, risk} \iqfirm{bank}]\label{iq:iv:cpp:10}
What does \cref{lst:iv:cpp:sign} print, and why is it a risk-control bug?
\end{interviewq}

\begin{interviewq}[$\star\star\star$ \iqroles{developer} \iqfirm{market maker}]\label{iq:iv:cpp:11}
Write the move constructor and move assignment operator of a class that owns a heap array of \texttt{int64\_t} and its size.
Why must they be \texttt{noexcept}? What state does a moved-from \texttt{std::vector} have?
\end{interviewq}

\begin{interviewq}[$\star\star\star$ \iqroles{developer} \iqfirm{proprietary firm}]\label{iq:iv:cpp:12}
Write a fixed-capacity, single-threaded ring buffer for a hot path. What happens when it is full, and why a power-of-two
capacity?
\end{interviewq}

\begin{interviewq}[$\star\star\star$ \iqroles{developer, mle} \iqfirm{market maker}]\label{iq:iv:cpp:13}
Write a compile-time table of the powers of ten that fit in an \texttt{int64\_t}, and a concept that accepts any type with
integer \texttt{price} and \texttt{qty} members. What happens if your table's loop computes $10^{19}$?
\end{interviewq}

\omcode{code/interviews/22-cpp/cpp/snippets/static_init.cpp}{1}{11}{Two global
variables.}\label{lst:iv:cpp:static}

\begin{interviewq}[$\star\star\star$ \iqroles{developer} \iqfirm{any}]\label{iq:iv:cpp:14}
What does \cref{lst:iv:cpp:static} print? What changes if \texttt{a} and \texttt{b} are defined in different source files,
and how do you make the program correct in that case?
\end{interviewq}

\begin{omsources}
\item ISO/IEC 14882:2020, \emph{Programming languages --- C++}, and the public working draft N4861 (2020), for the rules of
lifetime, overload resolution, constant evaluation and undefined behaviour.
\item One Quant Book~13, chapters 6--9 (C++ for latency; undefined behaviour and the compiler; Rust).
\end{omsources}
